\renewcommand{\Anno}{2013}

\renewcommand{\NumeroProblema}{1}

\renewcommand{\Testo}{
Per ogni numero reale $x$, indichiamo con $\lfloor x\rfloor$ 
la \emph{parte intera} di $x$, definita come il pi\`{u} grande 
intero minore od uguale ad $x$.

\begin{enumerate}
	\item[(a)]  Determinare per quali valori del parametro reale $a>0$ 
	l'equazione
	$$ax^{2}+10^{\lfloor x\rfloor}=2013$$
	ammette almeno una soluzione reale $x>0$.

	\item[(b)]  Determinare per quali valori del parametro reale $a>0$ 
	l'equazione
	$$ax^{2}+10^{-\lfloor x\rfloor}=2013$$
	ammette almeno una soluzione reale $x>0$.
\end{enumerate}
}

\renewcommand{\Risposta}{
\begin{enumerate}
	\renewcommand{\labelenumi}{(\alph{enumi})}
	\item  L'equazione ha almeno una soluzione $x>0$ se e solo se $a$ 
	verifica una delle seguenti 4 relazioni:
	$$\frac{1013}{16}<a\leq\frac{1013}{9},
	\quad\quad\quad
	\frac{1913}{9}<a\leq\frac{1913}{4},
	\quad\quad\quad
	\frac{2003}{4}<a\leq 2003,
	\quad\quad\quad
	2012<a.$$

	\item  L'equazione ha almeno una soluzione $x>0$ per ogni $a>0$. 
\end{enumerate}
}

\renewcommand{\Aiutino}{
Scrivere l'equazione portando la potenza di 10 a destra, quindi 
cercare di capire come sono fatti i grafici delle funzioni a destra e 
sinistra. A destra c'\`{e} una funzione a gradini, a sinistra una 
famiglia di parabole~\ldots

Questo permette di capire come vanno le cose. Per formalizzare il 
discorso algebricamente, conviene cercare di capire quando 
l'equazione ha una soluzione con $\lfloor x\rfloor=n$.}

\renewcommand{\SoluzioneUfficiale}{\TitoloSoluzione\DomandaA
Mostriamo che l'equazione ha almeno una soluzione $x>0$ se e solo se $a$ 
verifica una delle seguenti 4 relazioni:
$$\frac{1013}{16}<a\leq\frac{1013}{9},
\quad\quad\quad
\frac{1913}{9}<a\leq\frac{1913}{4},
\quad\quad\quad
\frac{2003}{4}<a\leq 2003,
\quad\quad\quad
2012<a.$$

Iniziamo osservando che, essendo $a>0$, per $x>0$ i due addendi a 
sinistra sono entrambi positivi, dunque in particolare $10^{\lfloor 
x\rfloor}\leq 2013$, da cui $0\leq\lfloor x\rfloor\leq 3$. Pertanto 
$\lfloor x\rfloor$ pu\`{o} assumere soltanto i valori 0, 1, 2, 3.
\begin{itemize}
	\item  Se $\lfloor x\rfloor=0$, allora l'equazione diventa 
	$ax^{2}=2013-10^{0}=2012$, la quale ha come soluzione positiva 
	$x=\sqrt{2012/a}$. Ora si deve avere che $0\leq x<1$, in modo da 
	avere $\lfloor x\rfloor=0$, e questo accade se e solo se
	$$0\leq\sqrt{\frac{2012}{a}}<1.$$
	
	La prima disuguaglianza \`{e} sempre verificata, la seconda vale 
	se e solo se $a>2012$.
	
	\item  Se $\lfloor x\rfloor=1$, allora l'equazione diventa 
	$ax^{2}=2013-10^{1}=2003$, la quale ha come soluzione positiva 
	$x=\sqrt{2003/a}$. Dovendo essere $\lfloor x\rfloor=1$, si 
	dovr\`{a} avere che $1\leq x<2$, il che \`{e} vero se e solo se
	$$1\leq\sqrt{\frac{2003}{a}}<2.$$
	
	Con facili calcoli si ottiene che questo accade se e solo se
	$$\frac{2003}{4}<a\leq 2003.$$

	\item  Se $\lfloor x\rfloor=2$, allora l'equazione diventa 
	$ax^{2}=2013-10^{2}=1913$, la quale ha come soluzione positiva 
	$x=\sqrt{1913/a}$. Dovendo essere $\lfloor x\rfloor=2$, si 
	dovr\`{a} avere che $2\leq x<3$, il che \`{e} vero se e solo se
	$$2\leq\sqrt{\frac{1913}{a}}<3.$$
	
	Con facili calcoli si ottiene che questo accade se e solo se
	$$\frac{1913}{9}<a\leq\frac{1913}{4}.$$

	\item  Se $\lfloor x\rfloor=3$, allora l'equazione diventa 
	$ax^{2}=2013-10^{3}=1013$, la quale ha come soluzione positiva 
	$x=\sqrt{1013/a}$. Dovendo essere $\lfloor x\rfloor=3$, si 
	dovr\`{a} avere che $3\leq x<4$, il che \`{e} vero se e solo se
	$$3\leq\sqrt{\frac{1913}{a}}<4.$$
	
	Con facili calcoli si ottiene che questo accade se e solo se
	$$\frac{1013}{16}<a\leq\frac{1013}{9}.$$
\end{itemize}

Abbiamo cos\`{i} ottenuto le 4 relazioni che descrivono l'insieme dei 
valori di $a$ per cui l'equazione ha almeno una soluzione $x>0$. La 
dimostrazione presentata mostra anche che tale soluzione positiva, 
per i valori di $a$ per cui esiste, \`{e} unica.

\DomandaB

Mostriamo che l'equazione ha almeno una soluzione $x>0$ per ogni 
$a>0$. Consideriamo l'insieme degli interi $k\geq 0$ tali che
\begin{equation}
	ak^{2}+10^{-k}\leq 2013.
	\label{defn:k0}
\end{equation}

Tale insieme \`{e} non vuoto in quanto contiene per lo meno $k=0$, ed 
ha solo un numero finito di elementi in quanto la (\ref{defn:k0}) non 
pu\`{o} valere se $k$ \`{e} sufficientemente grande (ad esempio per 
$k>\sqrt{2013/a}$). Sia quindi $k_{0}$ il pi\`{u} grande valore di 
$k$ che soddisfa la (\ref{defn:k0}). Si avr\`{a} dunque che
\begin{equation}
	ak_{0}^{2}+10^{-k_{0}}\leq 2013,
	\label{k0-1}
\end{equation}
\begin{equation}
	a(k_{0}+1)^{2}+10^{-(k_{0}+1)}> 2013.
	\label{k0-2}
\end{equation}

Mostriamo ora che l'equazione proposta ha una soluzione con $\lfloor 
x\rfloor=k_{0}$, che sar\`{a} quindi data dalla formula
$$x=\sqrt{\frac{2013-10^{-k_{0}}}{a}}.$$

Per avere che si tratta effettivamente di una soluzione, dobbiamo
verificare che $\lfloor x\rfloor=k_{0}$, cio\`{e}
$$k_{0}\leq\sqrt{\frac{2013-10^{-k_{0}}}{a}}< k_{0}+1.$$

Eleviamo al quadrato e moltiplichiamo per $a$. La prima disuguaglianza 
diventa 
$$ak_{0}^{2}\leq 2013-10^{-k_{0}},$$ 
che \`{e} equivalente alla 
(\ref{k0-1}). La seconda disuguaglianza diventa 
$$2013-10^{-k_{0}}<a(k_{0}+1)^{2},$$ 
la quale segue dalla (\ref{k0-2}) 
dal momento che 
$$2013-10^{-k_{0}}<2013-10^{-(k_{0}+1)}<a(k_{0}+1)^{2}.$$
}

\renewcommand{\SoluzioniAlternative}{

Presentiamo un approccio alternativo al punto~(b), pi\`{u} simile a
quello che abbiamo gi\`{a} adottato per il punto~(a).  Iniziamo a
chiederci per quali valori di $a$ esiste una soluzione con $\lfloor
x\rfloor=0$.  In tal caso l'equazione diventa $ax^{2}+1=2013$, da cui
$x=\sqrt{2012/a}$ (siamo infatti interessati solo alle soluzioni
positive).  Questa \`{e} effettivamente una soluzione se e solo se
$$0\leq\sqrt{\frac{2012}{a}}<1,$$
il che accade se e solo se $a>2012$.

Cerchiamo ora pi\`{u} in generale soluzioni con $\lfloor
x\rfloor=n$. Risolvendo ed imponendo la condizione $n\leq x<n+1$, 
troviamo che deve essere
$$n\leq\sqrt{\frac{2013-10^{-n}}{a}}<n+1,$$
da cui con facili calcoli troviamo le condizioni
$$\frac{2013-10^{-n}}{(n+1)^{2}}<a\leq\frac{2013-10^{-n}}{n^{2}}.$$

Questo mostra che, per ogni intero $n\geq 0$, l'equazione ha una
soluzione con $\lfloor x\rfloor=n$ se e solo se il parametro $a$
appartiene ad un opportuno insieme $I_{n}$.  Per $n=0$ tale insieme
\`{e} la semiretta $I_{0}=(2012,+\infty)$.  Per $n\geq 1$ l'insieme
$I_{n}$ \`{e} un intervallo del tipo $(A_{n},B_{n}]$ con
$$A_{n}=\frac{2013-10^{-n}}{(n+1)^{2}}, \hspace{3em}
B_{n}=\frac{2013-10^{-n}}{n^{2}}.$$

L'insieme degli $a$ per cui c'\`{e} soluzione \`{e} dunque l'unione di
$I_{n}$ al variare di $n\geq 0$.  Mostriamo ora che tale unione \`{e}
costituita da tutti gli $a>0$.  Questo segue da tre osservazioni.
Innanzitutto $I_{0}$ copre una semiretta.  Poi per ogni $n\geq 0$ si
ha che $$B_{n+1}=\frac{2013-10^{-(n+1)}}{(n+1)^{2}}>
\frac{2013-10^{-n}}{(n+1)^{2}}=A_{n},$$
quindi ogni intervallo $I_{n+1}$ si sovrappone al precedente $I_{n}$. 
Infine si ha che
$$\lim_{n\to +\infty}A_{n}=
\lim_{n\to +\infty}\frac{2013-10^{-n}}{(n+1)^{2}}=0,$$
per cui l'unione degli $I_{n}$ invade tutta la semiretta positiva.
}

\renewcommand{\Osservazioni}{Entrambi le equazioni ammettono una 
interpretazione grafica piuttosto semplice. Per il punto~(a), 
riscriviamo l'equazione nella forma
$$ax^{2}=2013-10^{\lfloor x\rfloor}.$$

Il grafico della funzione a destra, per $x>0$, \`{e} costituito da una
successione di ``scalini'' di lunghezza unitaria, di cui solo i primi
4 (corrispondenti a $\lfloor x\rfloor=0,1,2,3$) hanno altezza
positiva.  Il grafico della funzione $ax^{2}$ \`{e} una parabola
rivolta verso l'alto la cui ``apertura'' varia in funzione del
parametro $a$.

\`{E} chiaro quindi che per ciascuno dei 4 scalini esiste un 
intervallo di valori di $a$ per cui \`{e} intersecato dalla parabola. 
\`{E} chiaro anche che per $a$ sufficientemente grande la parabola 
\`{e} cos\`{i} stretta da intersecare sempre il primo scalino, quello 
corrispondente a $\lfloor x\rfloor=0$. Gli estremi degli intervalli 
di $a$ per cui c'\`{e} soluzione sono proprio i valori di $a$ per cui 
la parabola $ax^{2}$ passa per gli estremi degli scalini. 

La situazione \`{e} rappresentatata nella Figura~1a.  Gli unici 4
gradini che stanno nel primo quadrante hanno altezza 2012, 2003, 1913
e 1013, rispettivamente (per ovvie ragioni la figura non pu\`{o}
essere in scala).  Come si vede, esistono parabole che intersecano i
gradini, corrispondenti dunque a valori di $a$ per cui c'\`{e} una
soluzione $x>0$, necessariamente unica.  Ci sono per\`{o} anche
parabole, rappresentate tratteggiate, che passano tra un gradino e
l'altro senza attraversarne nessuno, corrispondenti quindi a valori di
$a$ per cui l'equazione non ha soluzioni positive.  \bigskip

{\centering
\hfill
\psset{unit=4ex}
\pspicture(-2,-3.5)(6,7)
\psline[linewidth=0.7\pslinewidth,linecolor=cyan]{->}(-1.5,0)(5.5,0)
\psline[linewidth=0.7\pslinewidth,linecolor=cyan]{->}(0,-2.5)(0,6.5)
\psline[linewidth=1.5\pslinewidth](-1,5.2)(0,5.2)
\psline[linewidth=1.5\pslinewidth](0,5)(1,5)
\psline[linewidth=1.5\pslinewidth](1,4.5)(2,4.5)
\psline[linewidth=1.5\pslinewidth](2,3.5)(3,3.5)
\psline[linewidth=1.5\pslinewidth](3,2)(4,2)
\psline[linewidth=1.5\pslinewidth](4,-2)(5,-2)
\psdots(-1,5.2)(0,5)(1,4.5)(2,3.5)(3,2)(4,-2)
\psclip{\psframe[linestyle=none](0,0)(6,6.5)}
\psset{linecolor=magenta}
\psplot{0}{5}{x x 10 mul mul}
\psplot{0}{5}{x x 2 mul mul}
\psplot[linestyle=dashed]{0}{5}{x x mul 3 div}
\psplot{0}{5}{x x 2 mul mul 9 div}
\psplot[linestyle=dashed]{0}{5}{x x mul 10 div}
\endpsclip
\rput(2,-3){Figura 1a}
\endpspicture
\hfill
\psset{unit=4ex}
\pspicture(-2,-3.5)(6,7)
\psline[linewidth=0.7\pslinewidth,linecolor=cyan]{->}(-1.5,0)(5.5,0)
\psline[linewidth=0.7\pslinewidth,linecolor=cyan]{->}(0,-2.5)(0,6.5)
\psline[linewidth=1.5\pslinewidth](-1,2)(0,2)
\psline[linewidth=1.5\pslinewidth](0,3)(1,3)
\psline[linewidth=1.5\pslinewidth](1,4)(2,4)
\psline[linewidth=1.5\pslinewidth](2,4.5)(3,4.5)
\psline[linewidth=1.5\pslinewidth](3,4.75)(4,4.75)
\psline[linewidth=1.5\pslinewidth](4,4.9)(5,4.9)
\psdots(-1,2)(0,3)(1,4)(2,4.5)(3,4.75)(4,4.9)
\psline[linecolor=orange](-1.5,5)(5.5,5)
\psclip{\psframe[linestyle=none](0,0)(6,6.5)}
\psset{linecolor=magenta}
\psplot{0}{5}{x x 10 mul mul}
\psplot{0}{5}{x x 2 mul mul}
\psplot{0}{5}{x x mul 4 div 4.25 mul}
\psplot{0}{5}{x x mul 3 div}
\psplot{0}{5}{x x 2 mul mul 9 div}
\endpsclip
\rput(2,-3){Figura 1b}
\endpspicture
\hfill
\mbox{}}

Per il punto~(b), riscriviamo l'equazione nella forma
$$ax^{2}=2013-10^{-\lfloor x\rfloor}.$$

La situazione \`{e} rappresentatata nella Figura~1b.  Ora ci sono
infiniti gradini nel primo quadrante, con altezze che partono da 2012
e crescono tendendo asintoticamente a 2013.  \`{E} evidente che tutte
le parabole intersecano almeno uno scalino, ma ci sono anche parabole,
come quella centrale nella figura, che intersecano 2 scalini.  Questo
mostra che ora la soluzione $x>0$ non \`{e} necessariamente unica per ogni
valore $a>0$ del parametro.  }

\renewcommand{\Commenti}{
}